\section{Results}
\label{sec:results}
We present our results in four parts, starting with shared prime factors
(Section~\ref{sec:shared-primes}), followed by shared moduli
(Section~\ref{sec:shared-moduli}), unusual exponents 
(Section~\ref{sec:unusual-exponents}), and finally targeted onion services
(Section~\ref{sec:targeted-onion-services}).

\subsection{Shared prime factors}
\label{sec:shared-primes}
Among all 588,945 identity keys, fastgcd found that 3,577 (0.61\%) 
moduli share prime factors.  We believed that 3,575 of these keys 
were all controlled by a single research group.  

Upon contacting the authors of the Security \& Privacy 2013 paper entitled
``Trawling for Tor hidden services''~\cite{Biryukov2013a}, we received
confirmation that these relays indeed were run by their research group and the
weak keys were caused by a shortcoming of their key generation utility. The
issue stemmed from the fact that their tool first generated thousands of prime
numbers and then computed multiple moduli using combinations of those prime
numbers in a greedy fashion without ensuring that the same primes were not
reused.

Because of the following shared properties, we are confident that all relays
were operated by the same group:

\begin{enumerate}
	\item All relays were online either in between November 11, 2012 and
		November 16, 2012 or in between January 14, 2013 and February 6, 2013,
		suggesting two separate experiments. We verified this by checking how
		long the relays stayed in the Tor consensus. The Tor consensus is
		updated hourly and documents which relays are available at a particular
		time. This data is archived by the Tor Project, and is made publically
		available on the CollectTor platform.

	\item All relays exhibited a predictable port assignment scheme.  In
		particular, we observed ports \{7003, 7007, \dots, 7043, 7047\} and
		\{8003, 8007, \dots, 8043, 8047\}.

	\item Except for two machines that were located in Russia and Luxembourg,
		all machines were hosted in Amazon's EC2 address space.  All machines
		except the one located in Luxembourg used Tor version 0.2.2.37.

	\item All physical machines had multiple fingerprints.  1,321 of these 3,557
		relays were previously characterized by Winter
		\ea~\cite[\S~5.1]{Winter2016a}.
\end{enumerate}

The remaining two keys belonged to a relay named ``DesasterBlaster,'' whose
origins we could not determine. Its router descriptor indicates that the relay
has been hosted on a MIPS machine which might suggest an embedded device with a
weak random number generator:

\begin{verbatim}
router DesasterBlaster 62.226.55.122 9001 0 0
platform Tor 0.2.2.13-alpha on Linux mips
\end{verbatim}

To further investigate, we checked whether the relay ``DesasterBlaster'' shares
prime factors with any other relays. It appears that the relay has rotated
multiple identity keys, and it only shares prime factors with its own keys.


\subsection{Shared moduli}
\label{sec:shared-moduli}
In addition to finding shared prime factors, we discovered relays that share a
\emph{modulus}, giving them the ability to calculate each other's private keys.
We know from ``Fact 1'' in~\cite[\S~1.1]{Boneh1999a} that two parties that share
a modulus $N$ can learn the prime factors $p$ and $q$ that produced their $N$.
With $p$, $q$, and each other's $e$'s in hand, the two parties can then compute
each other's decryption exponent $d$, at which point both parties now know the
public encryption and private decryptions keys. 

Table~\ref{tab:moduli} shows these ten relays with shared moduli clustered into
four groups. The table shows the relays' truncated, four-byte fingerprint, IP
addresses, and RSA exponents.  Note that the Tor client hard-codes the RSA
exponent to 65,537~\cite[\S~0.3]{torspec}, a recommended value that is resistant
to attacks against low public exponents~\cite[\S~4]{Boneh1999a}.  Any other
value indicates non-standard key generation.  All IP addresses were hosted by
OVH, a popular French hosting provider, and some of the IP addresses hosted two
relays, as our color coding indicates.  Finally, each group shared a four- or
five-digit prefix in their fingerprints.  We believe that a single attacker
controlled all these relays with the intention to manipulate the distributed
hash table that powers onion services~\cite{Biryukov2013a}---the shared
fingerprint prefix is an indication.  We suspect that the attackers iterated
over exponents to come up with the shared prefix.  The Tor Project informed us
that it discovered and blocked these relays in August 2014, when they first came
online~\cite{tor-priv1}.

\begin{table}[t]
	\centering
	\begin{tabular}{l l r}
	\toprule

	Short fingerprint & IP address & Exponent \\
	\midrule

	\texttt{\setlength{\fboxsep}{0pt}%
	\colorbox{highlight1}{\strut 838A}296A} & 188.165.164.163 &
	1,854,629 \\

	\texttt{\setlength{\fboxsep}{0pt}%
	\colorbox{highlight1}{\strut 838A}305F} &
	{\setlength{\fboxsep}{0pt}\colorbox{highlight3}{\strut 188.165.26.13}} &
	718,645 \\

	\texttt{\setlength{\fboxsep}{0pt}%
	\colorbox{highlight1}{\strut 838A}71E2} &
	{\setlength{\fboxsep}{0pt}\colorbox{highlight2}{\strut 178.32.143.175}} &
	220,955 \\

	\midrule

	\texttt{\setlength{\fboxsep}{0pt}%
	\colorbox{highlight1}{\strut 2249E}B42} &
	{\setlength{\fboxsep}{0pt}\colorbox{highlight3}{\strut 188.165.26.13}} &
	4,510,659 \\

	\texttt{\setlength{\fboxsep}{0pt}%
	\colorbox{highlight1}{\strut 2249E}C78} &
	{\setlength{\fboxsep}{0pt}\colorbox{highlight2}{\strut 178.32.143.175}} &
	1,074,365 \\

	\midrule

	\texttt{\setlength{\fboxsep}{0pt}%
	\colorbox{highlight1}{\strut E1EF}A388} & 188.165.3.63 &
	18,177 \\

	\texttt{\setlength{\fboxsep}{0pt}%
	\colorbox{highlight1}{\strut E1EF}8985} &
	{\setlength{\fboxsep}{0pt}\colorbox{highlight4}{\strut 188.165.138.181}} &
	546,019 \\

	\texttt{\setlength{\fboxsep}{0pt}%
	\colorbox{highlight1}{\strut E1EF}9EB8} &
	{\setlength{\fboxsep}{0pt}\colorbox{highlight5}{\strut 5.39.122.66}} &
	73,389 \\

	\midrule

	\texttt{\setlength{\fboxsep}{0pt}%
	\colorbox{highlight1}{\strut 410B}A17E} &
	{\setlength{\fboxsep}{0pt}\colorbox{highlight4}{\strut 188.165.138.181}} &
	1,979,465 \\

	\texttt{\setlength{\fboxsep}{0pt}%
	\colorbox{highlight1}{\strut 410B}B962} &
	{\setlength{\fboxsep}{0pt}\colorbox{highlight5}{\strut 5.39.122.66}} &
	341,785 \\

	\bottomrule
	\end{tabular}

	\caption{Four groups of relays that have a shared modulus.  All relays
	further share a fingerprint prefix in groups of two or three, presumably to
	manipulate Tor's distributed hash table.}
	\label{tab:moduli}
\end{table}

\subsection{Unusual exponents}
\label{sec:unusual-exponents}
The Tor source code hard-codes the public RSA exponent to 65,537, which is best
practice~\cite[\S~4]{Boneh1999a}.  Interested in non-standard key
generation, we checked if our data set featured relays with different
exponents.  Non-standard exponents suggest that a relay was after a specific
fingerprint in order to position itself in Tor's DHT.\footnote{A different
approach to detecting relays that position themselves in the DHT is to
determine how often they change their
fingerprints.~\cite[\S~4.3.3]{Winter2016a}} To obtain a fingerprint with a given
prefix, an adversary has to modify the underlying key material $p$, $q$, and $e$
until they result in the desired prefix.  Repeated modification of $e$ is
significantly more efficient than modifying $p$ or $q$ because it is costly to
verify if a large number is prime.  Leveraging this method, the tool
Scallion~\cite{scallion} generates vanity onion service domains by iterating
over the service's public exponent.

Among all of our 3.7 million keys, 122 possessed an exponent other than
65,537.\footnote{We list all relays in detail in
Appendix~\ref{sec:full-unusual-exponents}.} One relay had both non-standard
identity and onion key exponents while all remaining relays only had
non-standard identity key exponents.  Ten of these relays further had a shared
modulus, which we discuss in Section~\ref{sec:shared-moduli}.  Assuming that
these relays positioned themselves in the DHT to attack an onion service,
we were interested in learning what onion services they targeted.  It is
possible to identify the victims by compiling a comprehensive list of onion
services and determining each service's position in the DHT at the time
the malicious HSDirs were online.

\subsection{Identifying targeted onion services}
\label{sec:targeted-onion-services}

We then obtained a list of onion services by augmenting the list of the Ahmia
search engine~\cite{ahmia} with services that we discovered via Google searches
and by contacting researchers who have done similar work~\cite{Matic2015a}.  We
ended up with a list of 17,294 onion services, that were online at some point in
time.  Next, we developed a tool that takes as input our list of onion services
and the malicious HSDirs we discovered.\footnote{The tool is available online at
\url{https://github.com/citp/weak-tor-keys/tree/master/code/itos}.}  The tool
then determines all descriptors these onion services ever generated and checks
if any HSDir would have been close in the DHT to either of these
descriptors.

It is difficult to identify all targeted onion services with certainty.  First,
our list of onion services does not reveal when a service was online.  Second,
an HSDir could be responsible for an onion service simply by chance, instead of
on purpose, resulting in a false positive.  Third, our list of onion services is
not exhaustive, so we are bound to miss some potential victims.

Our tool identified the following four onion services for which we have strong
evidence that they were purposely targeted:

\begin{description}
	\item[\texttt{22u75kqyl666joi2.onion}] The service appears to be offline
		today, so we were unable to see for ourselves what it hosted.  According
		to cached index pages we found online, the onion service used to host a
		technology-focused forum in Chinese.  A subset of relays from
		Table~\ref{tab:group2} targeted the onion service on both August 14 and
		15, 2015 by providing nine out of the total of twelve responsible
		HSDirs.

	\item[\texttt{n3q7l52nfpm77vnf.onion}] As of February 2017, the service is
		still online, hosting the ``Marxists Internet Archive,'' an online
		archive of literature.\footnote{The onion service seems to be identical
		to the website \url{https://www.marxists.org}.} Figure~\ref{fig:archive}
		shows a screenshot of the service's index page.  A subset of relays from
		Table~\ref{tab:group1} targeted the onion service from November 27 to
		December 4, 2016.  The malicious HSDirs acted inconsistently,
		occasionally targeting only one replica.

	\item[\texttt{silkroadvb5piz3r.onion}] The onion service used to host the
		Silk Road marketplace, whose predominant use was a market for narcotics.
		The service was targeted by a subset of relays from
		Table~\ref{tab:group4}, from May 21 to June 3, 2013.  The HSDirs were
		part of a measurement experiment that resulted in a blog
		post~\cite{OCearbhaill2013a}.

	\item[\texttt{thehub7gqe43miyc.onion}] The onion service used to host a
		discussion forum, ``The Hub,'' focused on darknet markets.  A subset of
		relays from Table~\ref{tab:group2} targeted both of The Hub's replicas
		from August 22, 2015.
\end{description}

\begin{figure}[t]
	\centering
	\includegraphics[width=\linewidth]{figures/marxists-internet-archive.jpg}
	\caption{A screenshot of the index page of the onion service
		\texttt{n3q7l52nfpm77vnf.onion}, taken on February 13, 2017.}
	\label{fig:archive}
\end{figure}
